\documentclass[11pt]{article}
\usepackage[margin=1in]{geometry}
\usepackage{amsmath,amssymb}
\usepackage[T1]{fontenc}
\usepackage{microtype}
\setlength{\parskip}{3pt}
\newcommand{\Vb}{V}
\title{\vspace{-1.4em}\textbf{Design note: sequential fertility with a priced biological clock}\\[2pt]
\large A nested modification of the intergenerational housing--fertility model}
\author{Fertility--housing project}
\date{July 18, 2026 --- design for the \texttt{intergen\_seq\_fertility} fork; no production change}
\begin{document}
\maketitle
\vspace{-1.4em}

\section*{1. Purpose and scope}

In the production model, fertility is a \emph{one-shot family-size choice
with a recurring clock}: at each fertile age a childless household's
$\kappa_f$-logit picks among remaining childless and forming a family of
size $1,\dots,\bar n$; a chosen family arrives with certainty, lands in the
just-gave-birth stage, and the fertility choice never fires again. The
choice \emph{recurs} period by period only while childless. What the model
lacks is a \emph{price of time}: family formation succeeds with certainty
whenever chosen, so waiting is costless, the option value of postponement is
degenerate, and every response of fertility to risk, prices, or policy loads
on completed size rather than on timing. This note specifies the minimal
modification that prices time --- an age-declining conception probability
--- together with the new timing moments it makes measurable, the
identification bookkeeping, and the verification protocol. The goal of the
fork is a solver that is (i) bit-identical to production when the mechanism
is off, (ii) tested, and (iii) ready for calibration and for a
timing-versus-completed-fertility comparison against the production model.
The payoffs we are buying, beyond fit: policy experiments gain a
tempo/quantum decomposition (do birth grants create births or move them?),
housing affordability gains a postponement-induced childlessness channel,
and age-targeted policies (grants by age, entry-debt relief) acquire an
economic rather than mechanical age profile. Steady-state comparisons reveal
tempo in hazard shapes; full announcement-dynamics of tempo require
transition paths, which remain out of scope here.

Three margins are deliberately \emph{not} added in v1, each one line of why:
contraception failure (unwanted births) would require family-planning
moments we do not target; attempt costs (medical, search) are unidentified
without ART data; within-couple bargaining (Doepke--Kindermann) changes the
decision unit, which is a different paper.

\section*{2. The modification, as primitives}

\textbf{Current primitive.} During fertile ages $j \in [A_f^{start},
A_f^{end}]$ a childless household at $(n,s)=(0,0)$ faces a
$\kappa_f$-smoothed logit over $\{$remain childless$\} \cup \{$form a family
of size $n'\in\{1,\dots,\bar n\}\}$; the chosen family arrives with
certainty at $(n',1)$, child stages then age by the existing
stage-transition machinery, and households with $n\ge 1$ never face the
fertility choice again.

\textbf{Modified primitive.} The childless household chooses to
\emph{attempt} family formation at size $n'$. An attempt succeeds within the
period with probability $\pi_j$, landing at $(n',1)$ exactly as today;
with probability $1-\pi_j$ it fails, the household remains at $(0,0)$, and
may attempt again next period inside the window --- retrial is what the
existing recurring choice already supports for free. The attempt itself is
costless in v1. The schedule $\pi_j$ is a primitive of the environment,
\[
  \pi_j \;=\; \max\bigl\{\,0,\;\min\{1,\;1-\omega_1 e^{\omega_2\,(a_j -
  A_f^{start})}\}\bigr\},
  \qquad \pi_j = 0 \text{ for } a_j \ge 45,
\]
the four-year-period analogue of the exponential infertility risk in Sommer
(2016), with $(\omega_1,\omega_2)$ calibrated \emph{externally} to
demographic conception-within-window schedules by maternal age
(L\'eridon-type fecundity tables aggregated to the model's four-year
periods), not estimated in the SMM. Setting $\pi_j \equiv 1$ recovers the
production model identically; this nesting is the fork's acceptance test.

\textbf{Choice and value.} Each family-formation alternative in the
fertile-age logit replaces its certain continuation with the lottery
\[
  \Vb^{attempt\,n'}_j(b,z,\cdot) \;=\; \pi_j\,\Vb_j(b,z,n',1,\cdot)
  \;+\; (1-\pi_j)\,\Vb_j(b,z,0,0,\cdot),
\]
inside the same $\kappa_f$-logit over alternatives (the remain-childless
alternative is unchanged). One sentence of
economics: because $\pi_j$ falls with age, the value of postponing an
attempt now carries a real cost --- the risk that the window closes --- so
delay is an insurance margin with a price, and households near the end of
the window attempt more urgently than the same households earlier in life.
Nothing else in the household problem changes: budget sets, housing menus,
tenure choice, location choice, bequests, and the transfer floor added on
July 18 are untouched.

\textbf{Distribution.} Wherever the population flow currently moves
attempting mass into the birth state with certainty, it now splits it
$\pi_j : (1-\pi_j)$ between the birth destination and the unchanged state.
The accounting identity ``births into $(n+1)$ at age $j$ $=$ $\pi_j \times$
attempting mass at $j$'' must hold exactly and is a unit test. Measured
births everywhere in the model (birth grants, entrant creation from matured
children, completed fertility entering bequest utility) follow realized
births, not attempts, so all downstream objects update through the same
realized flow with no further modification.

\section*{3. New measurement}

All of the following are read off population flows; none require simulation.
Timing resolution is one model period (four years) --- stated once,
honestly: age at first birth is measured in four-year bins, and spacing in
whole periods.

\begin{itemize}
\itemsep1pt
\item \textbf{Family-formation hazard} $h_j$: realized family formations at
age $j$ divided by childless mass at $j$; the attempted hazard $h^A_j$
(attempts over childless mass) is also computable, and $h_j = \pi_j h^A_j$
by construction. Because the production choice is one-shot in size, v1 has
no second-birth hazard and no spacing distribution --- families of size two
form in a single event. Pricing \emph{per-birth} sequencing (a real
second-birth hazard) requires redesigning the child state to track the
youngest child's stage jointly with a running birth count; that is the
costed v2 extension of Section 5, not part of this fork.
\item \textbf{Age at first birth}: the distribution of realized
family-formation flows across age bins; headline moments are the mean and
the share of first births at ages 30 and above.
\item \textbf{Childlessness decomposition}: completed childlessness splits
into \emph{chosen} (never attempted during the window) and \emph{clock}
(attempted at least once, never succeeded). Failed attempters and
never-attempters occupy the same state, so the exact split requires
carrying a shadow never-attempted mass through the population flow. In v1
the split is instead computed by integrating the measured attempt and
realized hazards over a synthetic cohort --- exact under within-age
homogeneity of attempt rates and labeled as an approximation otherwise;
the exact shadow-mass accounting is a v1.1 refinement, and either way it is
measurement, not a state: the household problem is unchanged.
\item \textbf{Tempo/quantum readout for policy}: for any policy
counterfactual, report the change in completed fertility (quantum) and the
change in the birth-hazard age profile and mean age at birth (tempo). In
steady-state comparisons the two are linked by cohort accounting; the
decomposition becomes fully dynamic only on transition paths (out of
scope).
\end{itemize}

\section*{4. Identification bookkeeping}

New parameters: $(\omega_1,\omega_2)$ and the window close --- all
\emph{external} restrictions from demographic schedules; the SMM gains no
free parameters. New usable moments: the age profile of first-birth hazards, mean age at
first birth, the over-30 first-birth share, and the chosen/clock
childlessness split at 45 --- from NCHS/NSFG vital statistics and the PSID
childbirth history (targets and standard errors to be constructed in a
separate data pass; that pass is the remaining precondition for
calibration). These moments over-identify the existing fertility block
($\psi$, $\kappa_f$, $\bar c_n$, $\bar h_n$, $\bar h_{jump}$), whose
current discipline is levels and gradients only. One substitution risk to
watch: the logit scale $\kappa_f$ and the fecundity schedule both smooth
hazards; they are separated because $\pi_j$ is pinned externally by
biology, leaving $\kappa_f$ to the choice margin (income and price
gradients of hazards). A known fixed-$\theta$ tension to expect: the clock
adds a force for \emph{early} births (attempt while $\pi$ is high), while
the production calibration generates late first births through savings and
housing constraints; how these interact under recalibration is exactly the
timing-versus-completed-fertility question the fork exists to answer.

\section*{5. Verification protocol and deliverables}

(i) \emph{Nesting}: with $\pi_j \equiv 1$ the fork must reproduce the
production M5 local baseline bit for bit (the same gate0 standard used for
the July-18 transfer floor: loss and market residual to the last digit).
(ii) \emph{Unit tests}: fecundity-schedule arithmetic; the birth-flow
accounting identity; hazard identities $h_j=\pi_j h^A_j$; conservation of
mass across the $\pi:(1-\pi)$ split; a revive/kill pair mirroring the
floor-test pattern. (iii) \emph{Smoke}: one GE solve at the M5 winner with
a demographic $\pi$ schedule, reporting the standard packet plus the new
timing moments --- the first quantitative look at what the clock does at
fixed parameters. (iv) \emph{Ready for calibration}: the fork solves, the
timing moments exist as model output with a documented regeneration
command, and the comparison protocol is: calibrate the fork on the
production target system plus timing moments, then compare both models'
policy experiments with the tempo/quantum readout of Section~3.

The fork lives in \texttt{code/model/intergen\_seq\_fertility/} (a renamed
copy of the production package with its own tests and drivers); production
code is not modified. Design decisions above that later prove wrong are
cheap to revise --- the fork is disposable by construction.

\textbf{The v2 extension, costed.} True per-birth sequencing --- separate
attempt decisions for the first and second child, a second-birth hazard,
and a spacing distribution --- requires the child state to track (birth
count, stage of youngest child) rather than the current (family size, one
stage clock), touching the transition tensor, the value-function shape,
both Bellman implementations, three population-flow implementations, and
every $(n,s)$-keyed policy hook. It is the natural next step if the v1
comparison shows timing margins matter, and a separate scope decision: v1's
attempt machinery, fecundity schedule, timing moments, and tests all carry
over to it unchanged.

\end{document}
